% Paper 1, §3 — Tokenization, embedding, and unembedding
% Anchors: kb/notes/architecture/{tokenization,embeddings-and-tying}.md
%          kb/excerpts/{sennrich2016,kudo2018-sentencepiece,press2017-tied-embed,
%                       radford2019-gpt2,blt2024}.md

\section{Tokenization, embedding, and unembedding}
\label{sec:tokenization-embedding}

The discrete-token interface of a \glsterm{transformer}{Transformer} \glsterm{llm}{LLM} is implemented by three
parameter blocks that bracket the \glsterm{residual-stream}{residual stream}: a tokenizer that maps
text to integer IDs in $\{0, \dots, |V|-1\}$, an input embedding
$E_{\text{in}} \in \R^{|V| \times \mathsf{D}}$ that lifts those IDs into
the \glsterm{residual-stream}{residual stream} of shape $\BSH$, and an LM head
$E_{\text{out}} \in \R^{\mathsf{D} \times |V|}$ that projects the final
\glsterm{residual-stream}{residual stream} back to \glsterm{logits}{logits} of shape $\bsv$. The thesis of this
section is that the field has converged on a small lineage of choices
along this triple: subword \glsterm{bpe}{BPE} in either the \glsterm{sentencepiece}{SentencePiece} or
byte-level (tiktoken) family, with a clear scale-dependent split on
\glsterm{weight-tying}{weight tying} — tied at $\lesssim 1$B \glsterm{parameters}{parameters}, untied at $\gtrsim 8$B —
and a still-open contender, the tokenizer-free byte-patch model (BLT),
which is the first 2024 architecture to match a token-based 8B baseline
at fixed FLOPs \citep{blt2024}. The choices are independent, but they
interact: \glsterm{vocabulary}{vocabulary} size sets the embedding-table cost, which sets the
parameter share at which untying becomes free.

\subsection{Subword tokenization: BPE and SentencePiece}
\label{sec:tok-bpe}

\citet{sennrich2016} adapts \emph{byte-pair encoding} \citep[Gage
1994;][\S 3.2]{sennrich2016}\footnote{KB excerpt:
\texttt{kb/excerpts/sennrich2016.md\#sec-3-2}.} from data compression
to NLP word segmentation. The training procedure is greedy bottom-up:
initialize the symbol \glsterm{vocabulary}{vocabulary} with the character (or byte) alphabet,
represent each word as a sequence of those symbols plus an
end-of-word marker, then iterate the pair-merge loop. The reference
implementation as transcribed in \citet[\S 3.2, Algorithm 1]{sennrich2016}
is:

\begin{lstlisting}[language=Python, caption={BPE training (Sennrich et al.\ 2016, \S 3.2)}, label=lst:bpe]
import re, collections
def get_stats(vocab):
    pairs = collections.defaultdict(int)
    for word, freq in vocab.items():
        symbols = word.split()
        for i in range(len(symbols)-1):
            pairs[symbols[i], symbols[i+1]] += freq
    return pairs
def merge_vocab(pair, v_in):
    v_out = {}
    bigram = re.escape(' '.join(pair))
    p = re.compile(r'(?<!\S)' + bigram + r'(?!\S)')
    for word in v_in:
        w_out = p.sub(''.join(pair), word)
        v_out[w_out] = v_in[word]
    return v_out
num_merges = 10
for i in range(num_merges):
    pairs = get_stats(vocab)
    best = max(pairs, key=pairs.get)
    vocab = merge_vocab(best, vocab)
\end{lstlisting}

The final \glsterm{vocabulary}{vocabulary} size is the size of the initial alphabet plus the
number of merge operations — \emph{the only hyperparameter of the
algorithm} \citep[\S 3.2]{sennrich2016}. At inference, the same merge
list is replayed deterministically over new text. Crucially, \glsterm{bpe}{BPE} is
\emph{not} a left-to-right segmenter: the same prefix can yield a
different first token when the suffix changes. We will return to this
when contrasting with byte-patch models in \cref{sec:tok-blt}.

\citet{kudo2018-sentencepiece} packages \glsterm{bpe}{BPE} (and a unigram-LM
alternative) into a language-agnostic library with three additional
design choices that became the modern default for open-weight LLMs.
First, \emph{lossless \glsterm{tokenization}{tokenization}}: \glsterm{sentencepiece}{SentencePiece} treats raw input,
including whitespace, as a sequence of Unicode characters and escapes
spaces with the meta-symbol $\sqcup$ (U+2581) so that
$\mathrm{decode}(\mathrm{encode}(x)) = x$ holds verbatim
\citep[\S 3.1]{kudo2018-sentencepiece}.\footnote{KB excerpt:
\texttt{kb/excerpts/kudo2018-sentencepiece.md\#sec-3-1}.} Second, an
$O(N \log N)$ priority-queue training algorithm that scales to multi-TB
corpora \citep[\S 3.2]{kudo2018-sentencepiece}. Third, \emph{byte
fallback}: out-of-\glsterm{vocabulary}{vocabulary} characters decompose into the 256 UTF-8
byte tokens, eliminating the UNK symbol entirely
\citep[\S byte-fallback]{kudo2018-sentencepiece}. The three together
make a single tokenizer trainable on raw multilingual text — the choice
behind LLaMA-1/2/3 \citep{touvron2023, meta-llama3}, \glsterm{mistral}{Mistral}, Qwen,
Gemma, and OLMo 2 \citep{olmo2}.

\citet{radford2019-gpt2} addresses the orthogonal question of what
\emph{base alphabet} \glsterm{bpe}{BPE} should run on. A character-level base requires
the full Unicode code-point space ($> 130\text{K}$ symbols before any
merge), while a byte-level base needs only $256$. The naive
byte-level recipe, however, is sub-optimal for a load-bearing reason
\citep[\S 2.2]{radford2019-gpt2}\footnote{KB excerpt:
\texttt{kb/excerpts/radford2019-gpt2.md\#sec-2-2}.}:

\begin{quote}
\itshape We observed \glsterm{bpe}{BPE} including many versions of common words like
\texttt{dog} since they occur in many variations such as \texttt{dog.}
\texttt{dog!} \texttt{dog?}. This results in a sub-optimal allocation of
limited \glsterm{vocabulary}{vocabulary} slots and model capacity. To avoid this, we prevent
\glsterm{bpe}{BPE} from merging across character categories for any byte sequence. We
add an exception for spaces which significantly improves the compression
efficiency while adding only minimal fragmentation of words across
multiple vocab tokens.
\end{quote}

The category-respecting byte-level \glsterm{bpe}{BPE} of \citet{radford2019-gpt2} is
the direct ancestor of OpenAI's tiktoken (used by GPT-3, GPT-4, GPT-4o
at $|V|=200\text{K}$, and Phi-4 at $|V|=100\text{K}$
\citep[\S 2.1]{phi4}) and of \glsterm{precision-pinning}{DeepSeek-V3}'s $|V|=128\text{K}$ byte-level
\glsterm{bpe}{BPE} \citep[\S 2.1]{deepseek-v3}.

\subsection{Embedding and unembedding}
\label{sec:tok-embed-unembed}

Given token-ID sequence $\mathbf{t} \in \{0, \dots, |V|-1\}^N$, residual
width $\mathsf{D}$, and depth $L$, the input embedding is a learned
lookup,
\begin{equation}
X^0_i = E_{\text{in}}[t_i] \in \R^{\mathsf{D}}, \qquad
E_{\text{in}} \in \R^{|V| \times \mathsf{D}},
\label{eq:embed}
\end{equation}
producing the initial \glsterm{residual-stream}{residual stream} $X^0$ of shape $\BSH$. After $L$
blocks the LM head produces \glsterm{logits}{logits},
\begin{equation}
\mathrm{logits}_i = X^L_i \, E_{\text{out}} \in \R^{|V|}, \qquad
E_{\text{out}} \in \R^{\mathsf{D} \times |V|},
\label{eq:unembed}
\end{equation}
of shape $\bsv$ \citep[\S 3.4]{vaswani2017}. In modern \glsterm{decoder-only}{decoder-only}
LLMs $X^L$ is the output of the final \glsterm{rmsnorm}{RMSNorm} rather than the raw
post-block residual; \cref{eq:unembed} is otherwise unchanged. The
LM-head matmul costs $O(N \cdot |V| \cdot \mathsf{D})$ FLOPs per
position; at $|V| = 128\text{K}$ and $\mathsf{D} = 4096$ this is
$\sim\!5\times 10^{8}$ FLOPs per token, comparable to a single
\glsterm{transformer}{Transformer} block.

The original \glsterm{encoder-decoder}{encoder-decoder} \glsterm{transformer}{Transformer} scales the embedding by
$\sqrt{\mathsf{D}}$ before adding the positional signal
\citep[\S 3.4]{vaswani2017}, an artifact of the Xavier-initialization
variance regime; modern \glsterm{decoder-only}{decoder-only} LLMs typically omit the scaling
and rely on the first \glsterm{rmsnorm}{RMSNorm} to fix the norm of $X^0$.

\subsection{Weight tying}
\label{sec:tok-tying}

\citet{press2017-tied-embed} observed that the input embedding $U$ and
the output projection $W$ both index the same \glsterm{vocabulary}{vocabulary} and proposed
constraining them:
\begin{equation}
W = U, \quad \text{equivalently} \quad E_{\text{out}} = E_{\text{in}}^\top
\label{eq:tying}
\end{equation}
\citep[\S 3.1]{press2017-tied-embed}.\footnote{KB excerpt:
\texttt{kb/excerpts/press2017-tied-embed.md\#sec-3-1}.} The forward path
under tying is $\mathrm{logits}_i = X^L_i \, E_{\text{in}}^\top$, and a
single matrix $E_{\text{in}}$ is updated from both input and output
gradients. Tying halves the parameter count of the I/O blocks (saving
$|V| \cdot \mathsf{D}$ \glsterm{parameters}{parameters}) and the empirical result on \glsterm{rnn}{RNN}-era
language models was a \glsterm{perplexity}{perplexity} \emph{reduction}: PTB validation drops
from $78.4$ to $75.8$ on the medium \glsterm{lstm}{LSTM}, with similar gains on the
large \glsterm{lstm}{LSTM} and on Wikitext-2 \citep[\S 4]{press2017-tied-embed}. On NMT
\glsterm{encoder-decoder}{encoder-decoder} models with shared source/target \glsterm{vocabulary}{vocabulary}, three-way
tying reduced \glsterm{parameters}{parameters} by 28--52\% with no BLEU degradation
\citep[\S 5]{press2017-tied-embed}.

The frontier-2026 picture inverts the recommendation. \cref{tab:tying}
collects the public choices.

\begin{table}[h]
\centering
\small
\begin{tabular}{lllrr}
\toprule
Model & Tied? & $|V|$ & $\mathsf{D}$ & $|V|\mathsf{D}$ params \\
\midrule
Vaswani 2017 base & tied   & 37K   & 512   & 19M   \\
GPT-2 (1.5B)      & tied   & 50K   & 1600  & 80M   \\
LLaMA-1 / LLaMA-2 7B & tied   & 32K   & 4096  & 131M  \\
Phi-4             & tied   & 100K  & 5120  & 512M  \\
Qwen3 (small)     & tied   & 152K  & varies & varies \\
\midrule
LLaMA-3 8B        & untied & 128K  & 4096  & 525M $\times 2$ \\
LLaMA-3 70B/405B  & untied & 128K  & 8192/16384 & --- \\
OLMo 2 7B/13B     & untied & $\sim$50K & varies & varies \\
\glsterm{precision-pinning}{DeepSeek-V3} 671B  & untied & 128K  & 7168  & 920M $\times 2$ \\
Qwen3 (235B)      & untied & 152K  & varies & varies \\
\bottomrule
\end{tabular}
\caption{Tying decisions across reference models. Sources:
\citet[\S 3.4]{vaswani2017}, \citet[\S 2.3]{radford2019-gpt2},
\citet{touvron2023}, \citet[\S 2.1]{phi4},
\citet[\S 3.1]{meta-llama3}, \citet{olmo2},
\citet[\S 2.1]{deepseek-v3}, \citet{qwen3}.}
\label{tab:tying}
\end{table}

The empirical pattern is \emph{tied at $\lesssim 1$B, untied at
$\gtrsim 8$B}.

\begin{intuition}
Why frontier labs untie at scale: the embedding-block parameter saving
from tying is constant in $|V| \cdot \mathsf{D}$, while the
non-embedding model grows roughly as $L \cdot \mathsf{D}^2$. At LLaMA-3
8B the I/O embeddings are $\sim$12\% of model \glsterm{parameters}{parameters}; at 70B,
$\sim$5\%; at 405B, $\sim$1\%. The compression ratio of \cref{eq:tying}
goes from \emph{meaningful} to \emph{negligible} as scale grows, while
the alleged quality cost — the input embedding being pulled toward
output-space directions by the dominant LM-head gradient — is
constant. Returning to math: tying constrains
$E_{\text{out}}=E_{\text{in}}^\top$ globally, but the gradient
$\nabla_{E_{\text{in}}} \mathcal{L}$ acquires a contribution from the
\glsterm{softmax}{softmax} cross-entropy that scales with $|V|$, while the input-side
contribution scales only with the number of times each token is
embedded as input. The trade flips when $|V| \cdot \mathsf{D}$ ceases
to be a meaningful share of total \glsterm{parameters}{parameters}.
\end{intuition}

\begin{contradiction}
\citet[\S 4--5]{press2017-tied-embed} reports tying \emph{improves}
\glsterm{perplexity}{perplexity} on $\sim$10M-parameter \glsterm{lstm}{LSTM} language models. Frontier
2024--2026 open-weight LLMs at $\geq 8$B \glsterm{parameters}{parameters}
(\citealp{meta-llama3, olmo2, deepseek-v3, qwen3}) untie and report
this is the better choice. The reconciling claim — that the
tying-induced output-space bias only dominates once embeddings are a
small share of total \glsterm{parameters}{parameters} — is currently an
\emph{intuition}, not a derivation, and the few alternatives between
strict tying and full untying (e.g., pseudo-inverse tying,
$E_{\text{out}} = (E_{\text{in}}^\top)^+$) are recent and not yet
evaluated at frontier scale. \deepencite{press2017-tied-embed §6}
\end{contradiction}

\subsection{Vocabulary scaling}
\label{sec:tok-vocab}

\glsterm{vocabulary}{Vocabulary} size is the lever that connects tokenizer choice to
embedding cost. Smaller $|V|$ produces longer token sequences (higher
\emph{fertility}, defined as tokens-per-character), inflating
attention's $O(N^2)$ cost and reducing the model's effective context
horizon in characters. Larger $|V|$ produces shorter sequences but
inflates the embedding and LM-head matmul. The lineage of public
choices traces an upward arc: GPT-2 at $|V|=50{,}257$
\citep[\S 2.3]{radford2019-gpt2}; LLaMA-1/2 at $|V|=32{,}000$ on
\glsterm{sentencepiece}{SentencePiece} \glsterm{bpe}{BPE} \citep[\S 2.3]{touvron2023}; LLaMA-3 at
$|V|=128{,}256$ on expanded \glsterm{sentencepiece}{SentencePiece}-\glsterm{bpe}{BPE}
\citep[\S 3.1]{meta-llama3}; \glsterm{precision-pinning}{DeepSeek-V3} at $|V|=128{,}000$ on
byte-level \glsterm{bpe}{BPE} \citep[\S 2.1]{deepseek-v3}; Phi-4 at $|V|=100{,}000$
on tiktoken cl100k \citep[\S 2.1]{phi4}; GPT-4o at $|V|\!\approx\!200{,}000$
on tiktoken o200k; Qwen3 at $|V|=152{,}000$ \citep{qwen3}.

The capacity cost is concrete: LLaMA-3's $4\times$ vocab expansion
relative to LLaMA-2 added approximately $128{,}256 \times 4096 = 525\text{M}$
\glsterm{parameters}{parameters} per embedding matrix at 8B scale, and untied means double
that. The capability gain is also concrete: LLaMA-3 reports tokens-per-
English-word fertility near $1.0$ (vs.\ $\sim 1.5$ for LLaMA-2's 32K
vocab), with a substantially larger improvement on non-English text
where 32K \glsterm{sentencepiece}{SentencePiece}-\glsterm{bpe}{BPE} produced fertilities of $2$--$4$
\citep[\S 3.1]{meta-llama3}. The same expansion reserved a block of
unused token IDs explicitly for downstream fine-tuning, partly to
mitigate the \emph{glitch-token} failure mode where rarely-trained
embeddings sit near initialization and produce unpredictable model
behavior when prompted (the canonical example, the
\texttt{SolidGoldMagikarp} token in GPT-2/3, is a Tier C community
report and not relied on as a hard claim here).

\begin{analogy}
The tokenizer is a fixed prior the model cannot learn its way out of:
once $E_{\text{in}}$ is shaped by the merge list, every later
representation is conditioned on what the tokenizer chose to make a
single ID versus a sequence. If "12345" is one \glsterm{bpe}{BPE} merge but "1234" is
five characters, the model's digit-by-digit arithmetic competence is
not symmetric across strings — a fact LLaMA-3 addresses by tokenizing
digits individually \citep[\S 3.1]{meta-llama3}. Returning to math:
the tokenizer is the function $\mathbf{t}: \text{text} \to V^N$
upstream of \cref{eq:embed}; its merge list defines the equivalence
classes of strings that share a row of $E_{\text{in}}$, and the model
operates on those classes, not on the underlying bytes.
\end{analogy}

\subsection{Tokenizer-free: the BLT alternative}
\label{sec:tok-blt}

\citet{blt2024} replaces the fixed tokenizer with a learned
byte-patching scheme. A small byte-level auto-regressive language
model is trained alongside the main \glsterm{llm}{LLM} and used to score next-byte
entropy,
\begin{equation}
H(x_i) = \sum_{v \in V_{\text{byte}}} p_e(x_i = v \mid x_{<i}) \log p_e(x_i = v \mid x_{<i}),
\label{eq:blt-entropy}
\end{equation}
verbatim \citep[Eq.~1]{blt2024}.\footnote{KB excerpt:
\texttt{kb/excerpts/blt2024.md\#sec-2-3}.} Patch boundaries are placed
either when $H(x_t) > \theta_g$ (\emph{global} entropy threshold) or
when $H(x_t) - H(x_{t-1}) > \theta_r$ (\emph{approximate monotonic}
relative threshold) \citep[\S 2.3]{blt2024}. The result is an
\emph{incremental} patching function — the patching of $x_{<i}$ does
not depend on $x_{\geq i}$ — which \glsterm{bpe}{BPE} provably is not, because \glsterm{bpe}{BPE}
merges can re-segment a prefix when its continuation changes
\citep[\S 2.4]{blt2024}.

The architecture is three modules: a lightweight Local Encoder mapping
bytes to patch representations, a heavy Latent \glsterm{transformer}{Transformer} that runs
once per patch, and a lightweight Local Decoder that reads patch
outputs back to byte \glsterm{logits}{logits} \citep[\S 3]{blt2024}. The Latent
\glsterm{transformer}{Transformer}'s runtime scales with the number of patches $m$, not the
number of bytes $n$, and entropy patching produces $m \ll n$ on
predictable spans (e.g., the body of a named entity is one patch; the
boundary between sentences is its own patch). The headline scaling
result \citep[\S 1, Fig.~1]{blt2024}:

\begin{quote}
\itshape Overall, BLT matches training flop-controlled performance of
Llama 3 while using up to 50\% fewer flops at inference. … BLT models
unlock a new dimension for scaling LLMs, where model size can now be
scaled while maintaining a fixed-inference budget.
\end{quote}

\begin{forumsignal}
As of 2026-Q2, no production frontier-\glsterm{llm}{LLM} has shipped with a
BLT-family tokenizer-free architecture; the published 8B-scale parity
result \citep{blt2024} is the strongest public signal that
tokenizer-free is viable at non-trivial scale, and the FAIR follow-up
extending the comparison above 8B is reported in pre-print review but
not yet released. \deepencite{blt2024 §6}
\end{forumsignal}

\subsection{Forward link}

The choices in this section determine the rows of $E_{\text{in}}$ and
the columns of $E_{\text{out}}$, but they say nothing about \emph{order}
— two sequences that differ only by permutation produce the same
multiset of embedded vectors $\{E_{\text{in}}[t_i]\}_i$. The order
information is supplied by the positional-encoding choice that
follows: sinusoidal (\citealp[\S 3.5]{vaswani2017}), learned, \glsterm{t5}{T5}
relative bias, \glsterm{alibi}{ALiBi}, \glsterm{rope}{RoPE}, \glsterm{yarn}{YaRN}, \glsterm{nope}{NoPE}. The next section
(\cref{sec:positional}) walks that lineage, and explains why \glsterm{rope}{RoPE} plus
context-extension methods became the frontier-2026 default — a
choice that, like \glsterm{tokenization}{tokenization}, is one axis of the small set of
choices the modern \glsterm{transformer}{Transformer} actually offers.
